\section{Conclusion}\label{sec:conclusion}

In this work we initiated the study of EF orientations in multigraphs and presented a variety of both positive and negative results regarding
their computational complexity, including a complete characterization for the case of binary valuations. We additionally studied them under the perspective of parameterized complexity,
and showed how some of the results can even be transferred to the related {\EFXo} problem, improving upon previous work on the latter in the process and settling an open question by Deligkas et al.~\cite{ijcai/DeligkasEGK25}. Our work demonstrates that EF orientations are both interesting in their own right, and useful for transferring results to EFX orientations.

As a direction for future work, we comment that \cref{thm:EFo:Heavy} generalizes \cref{thm:EFo:linear_algo}, albeit only for simple graphs.
A natural question is thus whether we can develop such an algorithm that works for multigraphs as well, that is, whether {\EFo}
admits an FPT algorithm parameterized by the number $k$ of heavy edges running in time $2^k n^{\bO(1)}$.
Another interesting direction concerns our parameterized algorithms in \cref{subsec:EFo_shared_weight,subsec:EFXo:tw}, specifically, whether our parameterization by treewidth plus maximum shared weight can be improved to one by treewidth plus \emph{maximum weight}.

One last future research direction we would like to highlight is the search for regimes in which the complexity of EF and EFX \textsc{Orientation} differs.
We showed that one such case regards instances with binary valuations, while the two problems have similar behavior in the case of
polynomially-bounded weights and small vertex cover/treewidth.
In which other (valuation or structural) regime does the complexity of the two problems differ?



\subsubsection*{Acknowledgements}

We are grateful to our PhD advisors, Michael Lampis, Aris Pagourtzis, and Alkmini Sgouritsa, for valuable discussions and guidance regarding the problems studied in this work.

S.K., C.P., and M.S. were partially supported by project MIS 5154714 of the National Recovery and Resilience Plan Greece 2.0 funded by the European Union under the NextGenerationEU Program.


E.N. and M.V. were partially supported by the ANR project ANR-21-CE48-0022 (S-EX-AP-PE-AL).


S.K. and M.V. were partially supported by the project COAL-GAS,
financed under the Global Seed Fund program of Université Paris Sciences et Lettres (PSL) and by the CNRS.


S.K., C.P., and M.V. were partially supported by the Partenariat Hubert Curien France – Grèce 2025 programme ``Aristote'', funded by the Ministry of Europe and Foreign Affairs (MEAE) and the Ministry of Higher Education, Research, and Innovation (MESRI) of France, and the State Scholarships Foundation (IKY) of Greece, via the Cooperation and Cultural Action Service of the French Embassy, and the French Institute of Greece.

M.V. was partially supported by the ERC grant PARAMLP (project no. 101220563) funded by the European Union.
Views and opinions expressed are however those of the author(s) only and do not necessarily reflect those of the European Union or the European Research Council Executive Agency. Neither the European Union nor the granting authority can be held responsible for them.

S.K., C.P., and M.S. acknowledge the financial support of the Archimedes Research Unit (Athena RC) during their PhD.

No artificial intelligence was used at any point during the production of this paper. All of the presented results were conceived exclusively by the authors and the entirety of the paper was written exclusively by the authors.



%\manolis{Check these as well~\cite{ijcai/DeligkasEGHO22,ijcai/DeligkasEGHO21,ai/EibenGHO23,algorithmica/GanianOR23}.}


% THE FOLLOWING STEM FROM SECTION 5, COMMENTED THEM OUT

%From \cref{th:red_ef_alloc,thm:seth-ef} we obtain the following corollary.
%We conjecture that it is not tight; and explain it here

%\begin{corollary}
    %There is no $(2-\varepsilon)^kn^{\bO(1)}$ algorithm for \EFallo, where $k$ is the number of items with value at least $2$ for at least one agent.
%\end{corollary}
